\chapter{Elementary Data Structures}

So far in this book you have spent most of your energy on \emph{algorithms}: recipes that take some input and grind out an answer. Insertion sort, merge sort, quicksort, the order-statistic selection trick --- all of them assumed that the data was just \emph{sitting there}, in an array, ready to be poked at. We never asked the obvious follow-up question:

\begin{center}
\emph{How is the data actually stored, and what does that storage let us do quickly?}
\end{center}

That question is the subject of this chapter. A \textbf{data structure} is a specific way of organizing data in a computer's memory so that certain operations --- inserting a new item, deleting an old one, finding the item you want, asking ``what is next?'' --- become cheap. Choosing the right one is often the difference between a program that finishes in a blink and one that never finishes at all. The structures here are called \emph{elementary} not because they are unimportant, but because nearly every fancier structure later in the book (hash tables, balanced trees, heaps, graphs) is built out of these few primitives. Master these and you have the vocabulary for everything that follows.

We build, from the ground up: the humble \textbf{array} and its two-dimensional cousin the \textbf{matrix}; three disciplined ways to use an array --- the \textbf{stack}, the \textbf{queue}, and the \textbf{deque}; the pointer-and-object machinery that lets us escape the rigidity of arrays; \textbf{singly} and \textbf{doubly linked lists}; and how to represent \textbf{rooted trees}, which we lean on heavily for the rest of the book.

\begin{focusbox}
A data structure is the marriage of two things: a \emph{layout} (how bytes sit in memory) and a \emph{contract} (which operations it promises to support, and how fast). When you compare data structures, you are really comparing the costs of their contracts. Keep a running table in your head: for each structure, what does insert cost? Delete? Search? Access-by-position?
\end{focusbox}

\section{Memory, Addresses, and the Array}
\label{ch10:memory}

Before any data structure, there is memory. Picture your computer's memory as an enormously long street of identical houses, numbered $0, 1, 2, 3, \dots$ up into the billions. Each house holds one small fixed-size value (think of it as one byte, or for our purposes one ``cell''). The house number is called its \textbf{address}. The single most important fact about this street is:

\begin{intuitionbox}
Given an address, the hardware can fetch the value living there in \emph{constant time} --- one step, no matter how large the address. The machine does \emph{not} walk down the street counting houses. It teleports. This is why it is called \emph{random-access} memory: any house is as cheap to visit as any other.

Everything fast in this chapter is, at heart, an exploitation of this one teleportation power. Everything slow is what happens when you are forced to walk.
\end{intuitionbox}

An \textbf{array} is the most direct possible use of that power. You ask for a contiguous block of $n$ cells --- $n$ houses in a row --- and you get back the address of the first one, call it $base$. Now suppose each element of your array takes $s$ cells of space (an integer might be $s=4$ cells, say). Then element number $i$ lives at address
\[
  base + i \cdot s.
\]
That is a single multiplication and a single addition --- a constant amount of arithmetic --- followed by one teleporting memory fetch. So reading or writing $A[i]$ costs $\Th(1)$, \emph{regardless of which $i$ you ask for and regardless of how big the array is.}

\begin{focusbox}
The defining superpower of an array: \textbf{random access by index in $\Th(1)$ time}, via the address formula $base + i\cdot s$. This is why we have been able to write $A[i]$ in pseudocode all book long without worrying about cost. Now you know \emph{why} it is free.
\end{focusbox}

A concrete picture nails down the address arithmetic.

\begin{center}
\begin{tikzpicture}[font=\footnotesize]
  \def\w{1.0}
  \foreach \i/\val in {0/17, 1/3, 2/42, 3/8, 4/99} {
    \node[draw, minimum width=\w cm, minimum height=0.8cm, fill=blue!5] (c\i) at (\i*\w,0) {\val};
    \node[below=1pt of c\i, gray] {\small index \i};
  }
  \node[above=2pt of c0] {\small addr 1000};
  \node[above=2pt of c1] {\small 1004};
  \node[above=2pt of c2] {\small 1008};
  \node[above=2pt of c3] {\small 1012};
  \node[above=2pt of c4] {\small 1016};
  \draw[<->, sourceblue, thick] (-0.5,-1.1) -- (4.5,-1.1) node[midway, below] {one contiguous block};
\end{tikzpicture}
\end{center}

Here $base = 1000$ and each element takes $s = 4$ cells (a typical $32$-bit integer). The element $A[3] = 8$ lives at $1000 + 3 \cdot 4 = 1012$, exactly where the picture says. Notice index $3$ was no harder to reach than index $0$: both are one multiply, one add, one fetch.

\begin{example}[Addressing by hand]
\label{ch10:ex-addr}
An array of $64$-bit floating point numbers ($s = 8$ cells) starts at address $base = 200000$. Where does $A[10]$ live? Plug in:
\[
  200000 + 10 \cdot 8 = 200000 + 80 = 200080.
\]
And $A[11]$ is at $200088$, just $8$ cells further along. Consecutive elements sit $s$ cells apart --- that adjacency is the whole point of ``contiguous.''
\end{example}

\begin{trapbox}
\textbf{Zero-based versus one-based indexing.} CLRS and much of this book write arrays as $A[1..n]$, starting at $1$. Most real programming languages (C, Java, Python) start at $0$, so the first element is $A[0]$ and the last is $A[n-1]$. The address formula adapts simply: for one-based, element $i$ is at $base + (i-1)\cdot s$. The danger is the \emph{off-by-one} bug, the most common mistake in all of programming: writing $\le n$ where you meant $< n$, or forgetting to subtract $1$. Whenever you write a loop over an array, pause and check both endpoints explicitly: ``Does it touch the very first element? Does it touch the very last? Does it run one too far?''
\end{trapbox}

\subsection{The price of insertion and deletion}

Instant random access is powerful, but arrays have a hard limit baked into that contiguity. Suppose you have the array $[17, 3, 42, 8, 99]$ and you want to \emph{insert} the value $50$ so that it sits at index $2$ (pushing $42, 8, 99$ rightward). There is no spare house between index $1$ and index $2$; the block is packed solid. So you must \emph{shove} every element from index $2$ onward one house to the right to make a gap:

\begin{center}
\begin{tikzpicture}[font=\footnotesize]
  \def\w{1.0}
  % before
  \node at (-1.4,0.4) {before:};
  \foreach \i/\val in {0/17, 1/3, 2/42, 3/8, 4/99} {
    \node[draw, minimum width=\w cm, minimum height=0.7cm, fill=gray!8] (b\i) at (\i*\w,0.4) {\val};
  }
  % after
  \node at (-1.4,-0.9) {after:};
  \foreach \i/\val in {0/17, 1/3, 3/42, 4/8, 5/99} {
    \node[draw, minimum width=\w cm, minimum height=0.7cm, fill=gray!8] (a\i) at (\i*\w,-0.9) {\val};
  }
  \node[draw, minimum width=\w cm, minimum height=0.7cm, fill=green!18] at (2*\w,-0.9) {50};
  \draw[->, trapred, thick] (2*\w,0.05) to[bend left=15] (3*\w,-0.55);
  \draw[->, trapred, thick] (3*\w,0.05) to[bend left=15] (4*\w,-0.55);
  \draw[->, trapred, thick] (4*\w,0.05) to[bend left=15] (5*\w,-0.55);
\end{tikzpicture}
\end{center}

In the worst case --- inserting at the very front --- you have to shift all $n$ elements. So array insertion costs $\Th(n)$ in the worst case. Deletion is the mirror image: to remove the element at index $i$, you slide everything after it one house to the left to close the gap, again $\Th(n)$ in the worst case. The one happy exception is inserting or deleting at the \emph{end}: nothing needs to move, so it is $\Th(1)$ (provided there is room).

\begin{complexitybox}
\textbf{Array, costs.} Let $n$ be the number of elements.
\begin{center}
\begin{tabular}{lc}
\toprule
Operation & Worst-case time \\
\midrule
Access / update $A[i]$ & $\Th(1)$ \\
Search for a value (unsorted) & $\Th(n)$ \\
Insert / delete at the end & $\Th(1)$ \\
Insert / delete in the middle & $\Th(n)$ \\
\bottomrule
\end{tabular}
\end{center}
The array trades flexibility for raw access speed. It is the perfect choice when your data is fixed in size and you mostly read it, and a poor choice when items are constantly being inserted and removed in the middle.
\end{complexitybox}

\section{Matrices: arrays in two dimensions}
\label{ch10:matrix}

A \textbf{matrix} is a grid of values with $m$ rows and $n$ columns; we write $A[i,j]$ for the entry in row $i$, column $j$. But memory is a one-dimensional street --- there is no two-dimensional memory. So we must \emph{flatten} the grid into a single line of houses. There are two standard ways to do it, and they differ only in the order they visit the cells.

\textbf{Row-major order} lays down the entire first row, then the entire second row, and so on:
\[
  A[1,1],\, A[1,2],\, \dots,\, A[1,n],\ \ A[2,1],\, A[2,2],\, \dots
\]
\textbf{Column-major order} lays down the first column, then the second, and so on. C, Python (NumPy by default), and most languages use row-major; Fortran and MATLAB use column-major. The choice changes only the address formula. For an $m\times n$ matrix in row-major order with one-based indices, the entry $A[i,j]$ sits at offset
\[
  base + \big[(i-1)\cdot n + (j-1)\big]\cdot s.
\]

\begin{example}[Locating a matrix entry]
\label{ch10:ex-matrix}
Take a $3 \times 4$ matrix (so $m=3$, $n=4$), stored row-major with $s=1$ cell per entry and $base=0$. Where is $A[2,3]$? The full row $1$ uses indices $0,1,2,3$ (that is $n=4$ cells). Then within row $2$ we step to column $3$. The formula gives offset $(2-1)\cdot 4 + (3-1) = 4 + 2 = 6$. Lay out the flattened order to check:
\[
  \underbrace{A_{11},A_{12},A_{13},A_{14}}_{\text{offsets }0..3},\ \underbrace{A_{21},A_{22},\mathbf{A_{23}},A_{24}}_{\text{offsets }4..7},\ \dots
\]
Indeed $A[2,3]$ lands at offset $6$. The formula and the picture agree.
\end{example}

\begin{aestheticbox}
Row-major versus column-major is not academic trivia --- it can make real programs run several times faster or slower. Modern CPUs fetch memory in chunks (``cache lines''), so reading cells that are \emph{adjacent in memory} is far cheaper than jumping around. Loop over a row-major matrix row by row and you stride through memory in order, which the cache rewards. Loop column by column over that same matrix and you leap $n$ cells on every step, defeating the cache and slowing the loop dramatically. The math is identical; the running time is not. This is why numerical libraries care intensely about memory layout.
\end{aestheticbox}

\section{Stacks: last in, first out}
\label{ch10:stack}

Now we take an array and impose a \emph{discipline} on it. A \textbf{stack} is a collection where you may only add and remove from \emph{one end}, called the top. The rule is \textbf{LIFO}: last in, first out. Think of a stack of dinner plates: you add a clean plate to the top, you take the next plate from the top, and you never reach into the middle. The last plate placed is the first one taken.

\begin{intuitionbox}
A stack remembers things in reverse order of when they arrived. That sounds restrictive, but it is exactly the memory you want whenever you are working on something, get interrupted by a sub-task, and need to come back to where you were. Reading a book and following a footnote, which itself has a footnote? Your brain is running a stack. Pending function calls in a running program? A stack (literally called \emph{the} call stack). Matching brackets in an expression? A stack.
\end{intuitionbox}

A stack supports three operations:
\begin{itemize}
  \item \textsc{Push}$(S, x)$ --- put $x$ on top.
  \item \textsc{Pop}$(S)$ --- remove and return the top element.
  \item \textsc{Stack-Empty}$(S)$ --- is the stack empty?
\end{itemize}

Figure~\ref{ch10:fig-stack} illustrates the pointer mechanics. The array stays put; only $S.\mathit{top}$ moves.

\begin{figure}[H]
\centering
\begin{tikzpicture}[font=\small]
  \def\w{0.95}
  \foreach \i/\val in {1/4, 2/1, 3/3} {
    \node[draw, minimum width=\w cm, minimum height=0.75cm, fill=sourceblue!12] (s\i) at (\i*\w, 0) {\val};
  }
  \foreach \i in {4, 5} {
    \node[draw, minimum width=\w cm, minimum height=0.75cm, dashed, fill=gray!5] (s\i) at (\i*\w, 0) {};
  }
  \foreach \i in {1,...,5} {
    \node[below=2pt of s\i, gray, font=\scriptsize] {\i};
  }
  \node[left=3mm of s1] {$S:$};

  % Pointer S.top
  \node[sourceblue, font=\footnotesize, align=center] (toplab) at ($(s3.north)+(0, 1.05)$)
    {\textbf{$S.\mathit{top} = 3$}\\[-1pt]\scriptsize top element};
  \draw[->, sourceblue, very thick] (toplab.south) -- (s3.north);

  % Push annotation with generous clearance
  \draw[->, targetgreen, thick] ($(s4.north)+(0, 0.85)$) -- ($(s4.north)+(0, 0.1)$)
    node[midway, right=2pt, font=\scriptsize, targetgreen, align=left] {\textsc{Push} writes $S[4]$,\\[-2pt]$\mathit{top} \gets 4$};

  \draw[decorate, decoration={brace, amplitude=4pt}, gray] ($(s5.south east)+(0,-0.45)$) -- ($(s4.south west)+(0,-0.45)$)
    node[midway, below=4pt, gray, font=\scriptsize] {free capacity ($S.\mathit{size} = 5$)};
\end{tikzpicture}
\caption{A stack implemented on a fixed array $S[1\,..\,5]$. Live elements occupy slots $1$ through $S.\mathit{top} = 3$. \textsc{Push} increments $S.\mathit{top}$ and writes into the next free slot (green arrow); \textsc{Pop} reads $S[S.\mathit{top}]$ and decrements $S.\mathit{top}$. The array elements never shift, giving $\Th(1)$ worst-case time for all operations.}
\label{ch10:fig-stack}
\end{figure}

In Figure~\ref{ch10:fig-stack}, the stack holds $[4,1,3]$ with $top=3$; the next \textsc{Push} writes into slot $4$ and bumps $top$ to $4$, while the next \textsc{Pop} returns slot $3$ and drops $top$ to $2$. The dashed cells are reserved but unused.

\begin{algorithm}[H]
\caption{Stack operations on an array}
\DontPrintSemicolon
\SetKwProg{Fn}{}{}{}
\Fn{\textsc{Stack-Empty}$(S)$}{
  \lIf{$S.top == 0$}{\Return \textsc{true}}
  \lElse{\Return \textsc{false}}
}
\BlankLine
\Fn{\textsc{Push}$(S, x)$}{
  \If{$S.top == S.size$}{\Return \textsf{overflow} \tcp*{array is full}}
  $S.top \gets S.top + 1$\;
  $S[S.top] \gets x$\;
}
\BlankLine
\Fn{\textsc{Pop}$(S)$}{
  \If{\textsc{Stack-Empty}$(S)$}{\Return \textsf{underflow} \tcp*{nothing to pop}}
  $S.top \gets S.top - 1$\;
  \Return $S[S.top + 1]$\;
}
\end{algorithm}

Every line here is constant-time arithmetic and a single array access, so \textbf{push, pop, and empty-test all cost $\Th(1)$.} That is the whole appeal of a stack: its operations are about as cheap as operations get.

\begin{labbox}
Trace a stack by hand. Start empty and perform, in order: \textsc{Push}(4), \textsc{Push}(1), \textsc{Push}(3), \textsc{Pop}, \textsc{Push}(8), \textsc{Pop}, \textsc{Pop}.

We track the array contents and the value of $top$ after each step. Bold is the current top.
\begin{center}
\begin{tabular}{lll}
\toprule
Operation & Contents (bottom $\to$ top) & $top$ \\
\midrule
\textsc{Push}(4) & [\textbf{4}] & 1 \\
\textsc{Push}(1) & [4, \textbf{1}] & 2 \\
\textsc{Push}(3) & [4, 1, \textbf{3}] & 3 \\
\textsc{Pop} $\to 3$ & [4, \textbf{1}] & 2 \\
\textsc{Push}(8) & [4, 1, \textbf{8}] & 3 \\
\textsc{Pop} $\to 8$ & [4, \textbf{1}] & 2 \\
\textsc{Pop} $\to 1$ & [\textbf{4}] & 1 \\
\bottomrule
\end{tabular}
\end{center}
Notice the popped values come out $3, 8, 1$ --- exactly the reverse of the most recent pushes. Convince yourself that after the last step the stack still holds $4$ at the bottom, untouched since the very first push.
\end{labbox}

\begin{trapbox}
Two errors to guard against. \textbf{Underflow}: popping an empty stack. Always check \textsc{Stack-Empty} first, or you will read garbage at $S[0]$. \textbf{Overflow}: pushing onto a full fixed-size array. In the pseudocode above we return an error; in practice you would grow the array (more on that in the amortized-analysis chapter, where ``doubling'' makes the average push cost $\Th(1)$ even with growth). For now, just remember that a fixed array stack has a ceiling.
\end{trapbox}

\section{Queues: first in, first out}
\label{ch10:queue}

A \textbf{queue} flips the discipline around. You add at one end (the \textbf{tail}) and remove from the \emph{other} end (the \textbf{head}). The rule is \textbf{FIFO}: first in, first out --- exactly like a line of people waiting for coffee. The person who has waited longest is served next; newcomers join the back.

\begin{itemize}
  \item \textsc{Enqueue}$(Q, x)$ --- add $x$ at the tail.
  \item \textsc{Dequeue}$(Q)$ --- remove and return the element at the head.
\end{itemize}

The implementation is subtler than a stack, because the queue ``crawls'' down the array: every dequeue advances the head, every enqueue advances the tail, so the live region drifts rightward and would eventually fall off the end of a finite array. The fix is to make the array \textbf{circular} (or ``wrap around''): when an index reaches the end, it wraps back to the start, as if the array were a ring. We keep two indices, $Q.head$ (where the next dequeue happens) and $Q.tail$ (where the next enqueue will write). Figure~\ref{ch10:fig-circular-queue} illustrates the ring buffer and pointer movement.

\begin{figure}[H]
\centering
\begin{tikzpicture}[font=\small]
  \def\R{2.0}
  % Draw circular slots 0..7
  \foreach \i in {0,...,7} {
    \node[draw, circle, minimum size=0.88cm, fill=cyan!6] (q\i) at ({90 - \i*45}:\R) {};
    \node[gray, font=\scriptsize] at ({90 - \i*45}:\R+0.55) {slot \i};
  }
  % Active elements in slots 3..6
  \node[font=\bfseries] at (q3) {15};
  \node[font=\bfseries] at (q4) {6};
  \node[font=\bfseries] at (q5) {9};
  \node[font=\bfseries] at (q6) {8};

  % Clockwise motion indicator
  \draw[->, gray!60, line width=1.5pt] ({90 - 1.5*45}:0.95*\R)
    arc[start angle={90 - 1.5*45}, end angle={90 - 5.5*45}, radius=0.95*\R];
  \node[gray!70, font=\scriptsize, align=center] at (0, 0) {advance\\clockwise\\modulo $m$};

  % head pointer into slot 3 (bottom-right)
  \node[targetgreen, font=\footnotesize, align=left] (hl) at (3.8, -1.4)
    {\textbf{$\mathit{head} = 3$}\\[-1pt]\scriptsize next \textsc{Dequeue}\\[-2pt]\scriptsize advances $\to 4$};
  \draw[->, targetgreen, very thick] (hl.west) -- (q3.east);

  % tail pointer into slot 7 (top-left)
  \node[quotientorange, font=\footnotesize, align=right] (tl) at (-3.8, 1.4)
    {\textbf{$\mathit{tail} = 7$}\\[-1pt]\scriptsize next \textsc{Enqueue}\\[-2pt]\scriptsize wraps $\to 0$};
  \draw[->, quotientorange, very thick] (tl.east) -- (q7.west);

  % wrap arrow: index 7 -> index 0
  \draw[->, quotientorange, thick, dashed] ($(q7)+(0.45,-0.2)$) to[bend right=45]
    node[midway, right=2pt, font=\scriptsize\bfseries, quotientorange] {wrap $(7{+}1)\bmod 8 = 0$} ($(q0)+(0.45,0.2)$);
\end{tikzpicture}
\caption{The queue as a circular ring buffer ($m = 8$). Live elements ($15, 6, 9, 8$) occupy contiguous slots $3$ through $6$. \textsc{Dequeue} reads from $Q[\mathit{head}]$ and advances $\mathit{head}$ clockwise; \textsc{Enqueue} writes into $Q[\mathit{tail}]$ and advances $\mathit{tail}$. Modular arithmetic $(\mathit{tail} + 1) \bmod m$ wraps around from slot $7$ back to $0$ (dashed orange arrow), allowing the buffer to reuse space continuously in $\Th(1)$ time without shifting elements.}
\label{ch10:fig-circular-queue}
\end{figure}

In Figure~\ref{ch10:fig-circular-queue}, the live elements $15, 6, 9, 8$ sit at indices $3,4,5,6$; the next dequeue returns $15$ (at $head=3$), and the next enqueue writes at $tail=7$. The next index after $7$ wraps around to $0$. We do the wrap with modular arithmetic: ``the index after $i$'' in an array of size $m$ is $(i \bmod m) + 1$ in one-based terms, or simply $(i+1) \bmod m$ in zero-based terms.

\begin{algorithm}[H]
\caption{Queue operations on a circular array $Q[1..m]$}
\DontPrintSemicolon
\SetKwProg{Fn}{}{}{}
\Fn{\textsc{Enqueue}$(Q, x)$}{
  $Q[Q.tail] \gets x$\;
  \lIf{$Q.tail == Q.size$}{$Q.tail \gets 1$}
  \lElse{$Q.tail \gets Q.tail + 1$}
}
\BlankLine
\Fn{\textsc{Dequeue}$(Q)$}{
  $x \gets Q[Q.head]$\;
  \lIf{$Q.head == Q.size$}{$Q.head \gets 1$}
  \lElse{$Q.head \gets Q.head + 1$}
  \Return $x$\;
}
\end{algorithm}

Both operations are pure constant-time index arithmetic: \textbf{enqueue and dequeue cost $\Th(1)$.}

\begin{labbox}
Trace a circular queue of size $m=4$ (one-based indices $1..4$), starting with $head=1$, $tail=1$ (empty). Perform: \textsc{Enq}(7), \textsc{Enq}(2), \textsc{Enq}(5), \textsc{Deq}, \textsc{Enq}(9), \textsc{Enq}(1).

\begin{center}
\begin{tabular}{llll}
\toprule
Operation & Array (slots 1..4) & head & tail \\
\midrule
start & [\_,\_,\_,\_] & 1 & 1 \\
\textsc{Enq}(7) & [7,\_,\_,\_] & 1 & 2 \\
\textsc{Enq}(2) & [7,2,\_,\_] & 1 & 3 \\
\textsc{Enq}(5) & [7,2,5,\_] & 1 & 4 \\
\textsc{Deq} $\to 7$ & [7,2,5,\_] & 2 & 4 \\
\textsc{Enq}(9) & [7,2,5,9] & 2 & \textbf{1 (wrapped!)} \\
\textsc{Enq}(1) & [\textbf{1},2,5,9] & 2 & 2 \\
\bottomrule
\end{tabular}
\end{center}
Watch the wrap: after writing $9$ into slot $4$, $tail$ jumps back to $1$. The final \textsc{Enq}(1) overwrites slot $1$ --- which is legal because slot $1$'s old value ($7$) had already been dequeued. The live queue is now $2,5,9,1$ reading from head. Notice slot $1$ got reused without ever moving any data: that is what the circular layout buys you.
\end{labbox}

\begin{trapbox}
\textbf{Full versus empty ambiguity.} A subtle classic: in a circular array of size $m$, both a completely full queue and a completely empty queue can have $head == tail$. You cannot distinguish them from the indices alone! The standard fixes are (a) keep a separate count of live elements, or (b) never let the queue hold more than $m-1$ elements so $head == tail$ unambiguously means ``empty.'' Pick one and be consistent, or you will dequeue from an empty queue or silently overwrite live data.
\end{trapbox}

\section{Deques: open at both ends}
\label{ch10:deque}

A \textbf{deque} (pronounced ``deck,'' short for \emph{double-ended queue}) is the natural generalization: you may insert and delete at \emph{either} end. It subsumes both the stack and the queue --- use only one end and you have a stack; insert at one end and remove at the other and you have a queue. With a circular array and two indices it supports all four operations --- push-front, push-back, pop-front, pop-back --- in $\Th(1)$ time each. Deques return when we need a ``sliding window'' over a sequence; the point to remember is that open at both ends and still constant time is achievable.

\begin{walkbox}
Something to chew on while you walk. You are given only a stack (push/pop) as a black box, and you want to build a queue (enqueue/dequeue) out of it. A single stack reverses order, which is the \emph{opposite} of what a queue does --- so one stack alone cannot work. But \emph{two} stacks can. Can you figure out the trick? Hint: reversing a reversal gives you back the original order. What happens if you pour elements from one stack into another, one pop-then-push at a time? (This ``two stacks make a queue'' construction is a famous interview question, and it has a beautiful amortized cost of $\Th(1)$ per operation --- which you will be able to prove after the amortized-analysis chapter.)
\end{walkbox}

\section{Pointers and objects: escaping contiguity}
\label{ch10:pointers}

Every structure so far has leaned on the array's contiguous block. That block is exactly what makes random access fast --- and exactly what makes middle-insertion slow, because there is no slack between neighbors. What if we gave up on contiguity? What if each element could live \emph{anywhere} in memory, and simply held a note saying ``the next element is over there''?

That note is called a \textbf{pointer} (or reference): a value that is itself an address. An \textbf{object} is a small bundle of cells holding several named fields, one or more of which may be pointers to other objects. Given pointers, data structures stop being rigid blocks and become \emph{webs}: objects scattered across memory, wired together by pointers.

\begin{intuitionbox}
A pointer is an arrow. An object is a box with labeled slots, and some slots hold arrows pointing at other boxes. Drawing a data structure becomes literally drawing boxes and arrows. When you ``follow a pointer,'' you teleport (constant time!) to wherever it aims. The cost is that you can no longer compute where the $i$-th element lives by arithmetic --- there is no $base + i\cdot s$ formula across a web of scattered boxes. To reach the $i$-th element you must \emph{follow $i$ arrows}, one box at a time.
\end{intuitionbox}

We will write \lstinline|x.next| for ``the field named \emph{next} inside the object that $x$ points to,'' and we use the special value \textbf{\textsc{nil}} (sometimes written $\bot$ or \texttt{null}) for ``a pointer that points at nothing.'' \textsc{Nil} is how we mark the end of a chain or an empty structure --- it is the period at the end of the sentence.

\begin{trapbox}
The \textbf{null pointer} (\textsc{nil}) is the single most prolific source of crashes in real software --- so notorious its inventor, Tony Hoare (who also gave us quicksort), called it his ``billion-dollar mistake.'' The trap: trying to follow a pointer that is \textsc{nil}, i.e.\ asking for \lstinline|x.next| when $x$ itself is \textsc{nil}. There is no box there to look inside. Every single time you follow a pointer in a loop, you must first ask: ``Could this be \textsc{nil}?'' The loop condition \lstinline|while x != NIL| exists precisely to guard against this.
\end{trapbox}

\section{Linked lists}
\label{ch10:linked-list}

A \textbf{linked list} is the simplest pointer structure: a chain of objects, each called a \textbf{node}, where every node holds a \emph{key} (the actual data) and a pointer \emph{next} to the following node. The list as a whole is identified by a pointer \emph{head} to its first node. The last node's \emph{next} is \textsc{nil}, marking the end.

\begin{center}
\begin{tikzpicture}[font=\footnotesize,
   key/.style={draw, minimum width=0.7cm, minimum height=0.7cm, fill=structurepurple!8},
   nxt/.style={draw, minimum width=0.5cm, minimum height=0.7cm, fill=structurepurple!18}]
  \foreach \i/\val in {0/9, 1/16, 2/4, 3/1} {
    \node[key] (k\i) at (\i*2.0,0) {\val};
    \node[nxt] (n\i) at (\i*2.0+0.6,0) {};
  }
  \node[left=0.5cm of k0] (head) {\textit{head}};
  \draw[->, thick] (head) -- (k0.west);
  \draw[->, thick] (n0.east) -- (k1.west);
  \draw[->, thick] (n1.east) -- (k2.west);
  \draw[->, thick] (n2.east) -- (k3.west);
  \draw[->, thick] (n3.east) -- ++(0.5,0) node[right] {\textsc{nil}};
\end{tikzpicture}
\end{center}

This is a \textbf{singly linked list}: each arrow points only forward. There is also the \textbf{doubly linked list}, where each node additionally carries a \emph{prev} pointer to its predecessor, letting you walk backward as well as forward. Doubly linked lists make deletion much cleaner, as we will see, at the cost of one extra pointer per node.

\subsection{Searching a list}

How do you find a key in a linked list? There is no address arithmetic to jump to the middle --- you have no choice but to start at the head and follow \emph{next} pointers one node at a time until you find the key or hit \textsc{nil}.

\begin{algorithm}[H]
\caption{\textsc{List-Search}$(L, k)$ --- find the first node with key $k$}
\DontPrintSemicolon
\KwIn{a linked list $L$ and a key $k$}
\KwOut{a pointer to a node with key $k$, or \textsc{nil} if none exists}
$x \gets L.head$\;
\While{$x \neq \textsc{nil}$ \textbf{and} $x.key \neq k$}{
  $x \gets x.next$\;
}
\Return $x$\;
\end{algorithm}

In the worst case (the key is absent, or sits at the very end) you visit all $n$ nodes, so search costs $\Th(n)$ --- the same as an unsorted array, but \emph{slower in practice} because you follow scattered pointers instead of striding through contiguous cache-friendly memory. Linked lists earn their advantage elsewhere, in insertion and deletion.

\begin{formalbox}
\textbf{Loop invariant for \textsc{List-Search}.} \emph{At the start of each iteration of the while loop, none of the nodes strictly before $x$ in the list have key $k$.}

\emph{Initialization.} Before the first iteration $x = L.head$, so there are no nodes before $x$ at all; the claim holds vacuously.

\emph{Maintenance.} Suppose the invariant holds at the top of an iteration and the loop body runs. The body runs only when $x \neq \textsc{nil}$ and $x.key \neq k$, so the current node also fails to hold $k$. We then set $x \gets x.next$, advancing past the current node. Now \emph{every} node before the new $x$ --- those before the old $x$ (by hypothesis) plus the old $x$ itself (just checked) --- fails to hold $k$. The invariant is restored.

\emph{Termination.} The loop ends either because $x = \textsc{nil}$ (we walked off the end; by the invariant no earlier node held $k$, and there are no later ones, so $k$ is absent --- correctly returning \textsc{nil}) or because $x.key = k$ (we found it, and return that node). Either way the output is correct.
\end{formalbox}

\subsection{Inserting into a list}

Here is where linked lists shine. To insert a new node at the \emph{front} of the list, you do not shift anything --- you rewire two pointers. Make the new node point at the old head, then make $head$ point at the new node:

\begin{algorithm}[H]
\caption{\textsc{List-Prepend}$(L, x)$ --- insert node $x$ at the front}
\DontPrintSemicolon
\KwIn{a linked list $L$ and a fresh node $x$ to insert}
$x.next \gets L.head$\;
$L.head \gets x$\;
\end{algorithm}

Two pointer assignments, no shifting, no matter how long the list is: \textbf{front insertion is $\Th(1)$.} Contrast that with an array, where front insertion was $\Th(n)$ because everything had to slide. This is the linked list's whole reason to exist.

\begin{center}
\begin{tikzpicture}[font=\footnotesize,
   key/.style={draw, minimum width=0.7cm, minimum height=0.6cm, fill=structurepurple!8},
   nxt/.style={draw, minimum width=0.5cm, minimum height=0.6cm, fill=structurepurple!18}]
  \node[key] (ka) at (0,0) {16}; \node[nxt] (na) at (0.6,0) {};
  \node[key] (kb) at (2.0,0) {4}; \node[nxt] (nb) at (2.6,0) {};
  \node[key, fill=green!18] (kx) at (-0.5,1.4) {9}; \node[nxt, fill=green!30] (nx) at (0.1,1.4) {};
  \node[left=1.0cm of ka] (head) {\textit{head}};
  \draw[->, thick, green!45!black] (head) to[bend left=20] (kx.west);
  \draw[->, thick, green!45!black] (nx.south) to[bend left=10] (ka.north);
  \draw[->, thick] (na.east) -- (kb.west);
  \draw[->, thick] (nb.east) -- ++(0.45,0) node[right] {\textsc{nil}};
  \draw[->, thick, dashed, gray] (head) -- (ka.west) node[midway, below, gray] {(old)};
\end{tikzpicture}
\end{center}

The dashed gray arrow shows where $head$ used to point; the two green arrows are the rewiring. Nothing else in the list moved a single byte.

\subsection{Deleting from a list}

To delete a node, we splice it out: make its predecessor's \emph{next} skip over it and point straight at its successor (and, in a doubly linked list, make the successor's \emph{prev} point back at the predecessor). In a \textbf{doubly} linked list this is wonderfully direct, because every node knows its own predecessor through \emph{prev} --- no searching required.

\begin{center}
\begin{tikzpicture}[font=\footnotesize,
   nd/.style={draw, rounded corners=2pt, minimum width=0.9cm, minimum height=0.7cm}]
  \node[nd, fill=structurepurple!8] (p) at (0,0) {$a$};
  \node[nd, fill=trapred!12] (x) at (2.4,0) {$x$};
  \node[nd, fill=structurepurple!8] (s) at (4.8,0) {$b$};
  % original forward/back pointers (gray) through x
  \draw[->, gray] ($(p.north east)+(0,-0.08)$) to[bend left=22] ($(x.north west)+(0,-0.08)$);
  \draw[->, gray] ($(x.south west)+(0,0.08)$) to[bend left=22] ($(p.south east)+(0,0.08)$);
  \draw[->, gray] ($(x.north east)+(0,-0.08)$) to[bend left=22] ($(s.north west)+(0,-0.08)$);
  \draw[->, gray] ($(s.south west)+(0,0.08)$) to[bend left=22] ($(x.south east)+(0,0.08)$);
  % new splice arrows (green) directly between a and b
  \draw[->, targetgreen, very thick] ($(p.north east)+(0,0.18)$) to[bend left=30] ($(s.north west)+(0,0.18)$);
  \draw[->, targetgreen, very thick] ($(s.south west)+(0,-0.18)$) to[bend left=30] ($(p.south east)+(0,-0.18)$);
  \node[trapred, font=\scriptsize] at (2.4,-1.0) {(removed)};
  \node[targetgreen, font=\scriptsize] at (2.4,1.15) {$a.next \gets b,\ \ b.prev \gets a$};
\end{tikzpicture}
\end{center}

The gray arrows are the original links through $x$; the two bold green arrows are the new links that bypass $x$. Once nothing points at $x$, it is gone from the list.

\begin{algorithm}[H]
\caption{\textsc{List-Delete}$(L, x)$ --- remove node $x$ from a doubly linked list}
\DontPrintSemicolon
\KwIn{a doubly linked list $L$ and a pointer to the node $x$ to remove}
\lIf{$x.prev \neq \textsc{nil}$}{$x.prev.next \gets x.next$}
\lElse{$L.head \gets x.next$ \tcp*{$x$ was the head}}
\If{$x.next \neq \textsc{nil}$}{$x.next.prev \gets x.prev$}
\end{algorithm}

Given the pointer $x$ to the doomed node, deletion is a handful of pointer rewires --- \textbf{$\Th(1)$ time}, again independent of list length. (The \textsc{if}/\textsc{else} branches just handle the boundary cases where $x$ is the first node, which has no predecessor, or the last, which has no successor.) Compare the array, where deletion in the middle forced an $\Th(n)$ slide.

\begin{trapbox}
That last paragraph hid an assumption: ``\emph{given the pointer $x$}.'' If you only know the \emph{key} you want to delete, you must first \textsc{List-Search} for it --- which is $\Th(n)$ --- and \emph{then} splice it out in $\Th(1)$. So delete-by-key is $\Th(n)$ overall; only delete-by-pointer is truly $\Th(1)$. People love to claim ``linked list deletion is $O(1)$'' and forget the fine print. In a \emph{singly} linked list it is worse still: even with a pointer to $x$ you cannot reach $x.prev$ to rewire it, so you must walk from the head to find the predecessor --- $\Th(n)$ again. This single fact is the best argument for paying the extra pointer per node and going doubly linked.
\end{trapbox}

\subsection{The sentinel trick}

All those \textsc{nil}-checking \textsc{if}/\textsc{else} branches in \textsc{List-Delete} are ugly and bug-prone. A beautiful simplification is the \textbf{sentinel}: a single dummy node, $L.\mathit{nil}$, that has no real key but sits permanently in the list as both the node \emph{after} the last real node and the node \emph{before} the first. The list becomes \emph{circular}, with the sentinel stitching the two ends together. Now there are no boundaries: every real node always has a genuine predecessor and successor (possibly the sentinel itself), so the special cases vanish and \textsc{List-Delete} collapses to just two lines:
\[
  x.\mathit{prev}.\mathit{next} \gets x.\mathit{next}; \qquad x.\mathit{next}.\mathit{prev} \gets x.\mathit{prev}.
\]
Figure~\ref{ch10:fig-sentinel} shows the architecture: the circular sentinel structure eliminates boundary checks because every node, including the head and tail, always has valid neighbors.

\begin{figure}[H]
\centering
\begin{tikzpicture}[font=\small,
   nd/.style={draw, rounded corners=3pt, minimum width=1.1cm, minimum height=0.8cm},
   sentinel/.style={draw=structurepurple!80, line width=1.2pt, rounded corners=3pt, minimum width=1.3cm, minimum height=0.8cm, fill=structurepurple!15}]

  % Sentinel node at the left
  \node[sentinel] (nil) at (0, 0) {$L.\mathit{nil}$};

  % Real data nodes
  \node[nd, fill=sourceblue!12] (n1) at (2.4, 0) {$9$};
  \node[nd, fill=trapred!15] (n2) at (4.8, 0) {$16$};
  \node[nd, fill=sourceblue!12] (n3) at (7.2, 0) {$4$};

  % Node labels below
  \node[below=2pt of nil, font=\scriptsize, structurepurple] {sentinel};
  \node[below=2pt of n1, font=\scriptsize, gray] {head};
  \node[below=2pt of n2, font=\scriptsize, trapred] {doomed $x$};
  \node[below=2pt of n3, font=\scriptsize, gray] {tail};

  % Forward next pointers (top curve)
  \draw[->, thick, structurepurple] ($(nil.north east)+(0,-0.05)$) to[bend left=25]
    node[midway, above=1pt, font=\tiny, structurepurple] {\textit{next}} ($(n1.north west)+(0,-0.05)$);
  \draw[->, gray] ($(n1.north east)+(0,-0.05)$) to[bend left=25] ($(n2.north west)+(0,-0.05)$);
  \draw[->, gray] ($(n2.north east)+(0,-0.05)$) to[bend left=25] ($(n3.north west)+(0,-0.05)$);

  % Backward prev pointers (bottom curve)
  \draw[->, gray] ($(n3.south west)+(0,0.05)$) to[bend left=25] ($(n2.south east)+(0,0.05)$);
  \draw[->, gray] ($(n2.south west)+(0,0.05)$) to[bend left=25] ($(n1.south east)+(0,0.05)$);
  \draw[->, thick, structurepurple] ($(n1.south west)+(0,0.05)$) to[bend left=25]
    node[midway, below=1pt, font=\tiny, structurepurple] {\textit{prev}} ($(nil.south east)+(0,0.05)$);

  % Circular wrap links to/from sentinel
  \draw[->, thick, structurepurple] (n3.north) to[bend right=28]
    node[midway, above=2pt, font=\scriptsize, structurepurple] {tail's \textit{next} $\to L.\mathit{nil}$} ($(nil.north)+(0,0.1)$);
  \draw[->, thick, structurepurple] ($(nil.south)+(0,-0.1)$) to[bend right=28]
    node[midway, below=2pt, font=\scriptsize, structurepurple] {$L.\mathit{nil}.\mathit{prev} \to$ tail} (n3.south);

  % The two-line green splice bypassing x (n2)
  \draw[->, targetgreen, very thick] ($(n1.north east)+(0,0.15)$) to[bend left=32]
    node[midway, above=2pt, font=\scriptsize\bfseries, targetgreen] {$x.\mathit{prev}.\mathit{next} \gets x.\mathit{next}$} ($(n3.north west)+(0,0.15)$);
  \draw[->, targetgreen, very thick] ($(n3.south west)+(0,-0.15)$) to[bend left=32]
    node[midway, below=2pt, font=\scriptsize\bfseries, targetgreen] {$x.\mathit{next}.\mathit{prev} \gets x.\mathit{prev}$} ($(n1.south east)+(0,-0.15)$);

\end{tikzpicture}
\caption{The sentinel trick in a circular doubly linked list. The dummy node $L.\mathit{nil}$ stitches the head and tail together into a seamless ring: $L.\mathit{nil}.\mathit{next}$ points to the head, and $L.\mathit{nil}.\mathit{prev}$ points to the tail. Deleting any node $x$ --- whether in the middle, at the head, or at the tail --- executes the exact same two pointer assignments ($x.\mathit{prev}.\mathit{next} \gets x.\mathit{next}$ and $x.\mathit{next}.\mathit{prev} \gets x.\mathit{prev}$, bold green arrows), completely eliminating null-pointer checks and edge-case branching.}
\label{ch10:fig-sentinel}
\end{figure}

The cost is one wasted node's worth of memory and a little care during initialization. For lists where you do many insertions and deletions, the cleaner, branch-free code is well worth it.

\begin{complexitybox}
\textbf{Doubly linked list, costs.} Let $n$ be the number of nodes.
\begin{center}
\begin{tabular}{lc}
\toprule
Operation & Worst-case time \\
\midrule
Access $i$-th element & $\Th(n)$ \\
Search for a key & $\Th(n)$ \\
Insert at front / back$^\dagger$ & $\Th(1)$ \\
Delete a node (given its pointer) & $\Th(1)$ \\
Delete by key & $\Th(n)$ \\
\bottomrule
\end{tabular}
\end{center}
$^\dagger$Back insertion is $\Th(1)$ if you keep a $tail$ pointer; otherwise you must walk to the end, $\Th(n)$. The list is the array's complement: it gives up $\Th(1)$ random access to gain $\Th(1)$ structural edits.
\end{complexitybox}

\begin{aestheticbox}
When should you actually reach for a linked list over an array? When your workload is dominated by inserting and deleting at known positions and you rarely need to jump to the $i$-th element. The textbook real-world example is a music or video \emph{playlist} with constant reordering, or the list of free memory blocks an operating system splices in and out as programs allocate and free memory. But heed a modern caveat: because array elements sit contiguously and linked-list nodes are scattered, arrays are so much friendlier to the CPU cache that in practice a simple array often beats a linked list \emph{even for tasks the list is supposedly better at}, until $n$ gets quite large. Asymptotics tell you the trend; they do not always tell you the winner at the sizes you care about. Measure.
\end{aestheticbox}

\section{Representing rooted trees}
\label{ch10:trees}

We close with the structure that dominates the rest of the book: the \textbf{rooted tree}. A tree is a collection of nodes connected so that there is exactly one special node, the \textbf{root}, at the top, and every other node has exactly one \emph{parent} (the node just above it) and any number of \emph{children} (the nodes just below). No node is its own ancestor; there are no cycles. A node with no children is a \textbf{leaf}.

The pointer-and-object idea generalizes perfectly: a tree node is just an object holding a key plus pointers to its relatives. The only question is \emph{which} relatives to store pointers to, and that choice depends on how many children a node may have.

\subsection{Binary trees}

The cleanest case is the \textbf{binary tree}, where every node has at most two children, named \emph{left} and \emph{right}. Each node $x$ stores three pointers: $x.p$ to its parent, $x.left$ to its left child, $x.right$ to its right child. Any missing relative is \textsc{nil} --- the root's parent is \textsc{nil}, and a leaf's children are \textsc{nil}.

\begin{center}
\begin{tikzpicture}[font=\footnotesize, level distance=1.1cm,
   level 1/.style={sibling distance=3cm},
   level 2/.style={sibling distance=1.5cm},
   every node/.style={draw, circle, minimum size=0.7cm, fill=targetgreen!12}]
  \node {18}
    child {node {7}
      child {node {3}}
      child {node {12}}
    }
    child {node {26}
      child[missing]
      child {node {31}}
    };
\end{tikzpicture}
\end{center}

Here $18$ is the root. Its left child is $7$, its right child is $26$. Node $7$ has two children, $3$ and $12$ (both leaves). Node $26$ has \emph{no} left child (that pointer is \textsc{nil}) and one right child, $31$. Reading the tree from the top, each node points down to up to two children and up to one parent.

\subsection{When the number of children is unbounded}

Binary trees are tidy because ``at most two children'' means a fixed two pointer slots. But many trees let a node have \emph{arbitrarily} many children --- think of a file system folder that may contain hundreds of items, or a company org chart. We cannot give each node a fixed pointer for ``child 1, child 2, child 3, \dots'' because we do not know in advance how many there will be, and reserving a huge fixed number would waste enormous memory on nodes with few children.

The space-frugal answer is the \textbf{left-child, right-sibling} representation. Each node stores only \emph{two} pointers, regardless of how many children it has:
\begin{itemize}
  \item \emph{left-child}: a pointer to its \emph{first} child (or \textsc{nil} if it is a leaf);
  \item \emph{right-sibling}: a pointer to its \emph{next sibling} --- the next child of the same parent (or \textsc{nil} if it is the last child).
\end{itemize}

\begin{intuitionbox}
Picture each parent reaching down with a single arrow to its \emph{eldest} child. The siblings then hold hands left to right in a chain. To visit \emph{all} the children of a node, you drop down to the first child, then walk the sibling chain to the right until it runs out. You have turned ``a list of children of unknown length'' into exactly the linked list from the previous section --- so it costs only two pointers per node, no matter how lopsided or bushy the tree.
\end{intuitionbox}

\begin{center}
\begin{tikzpicture}[font=\footnotesize,
   nd/.style={draw, circle, minimum size=0.7cm, fill=structurepurple!12}]
  \node[nd] (r) at (0,0) {A};
  \node[nd] (b) at (-1.5,-1.4) {B};
  \node[nd] (c) at (0,-1.4) {C};
  \node[nd] (d) at (1.5,-1.4) {D};
  \node[nd] (e) at (-2.2,-2.8) {E};
  \node[nd] (f) at (-0.8,-2.8) {F};
  % left-child arrows (solid, down to first child)
  \draw[->, thick, structurepurple] (r.south) to[bend right=10] (b.north);
  \draw[->, thick, structurepurple] (b.south) to[bend right=10] (e.north);
  % right-sibling arrows (dashed, across)
  \draw[->, thick, dashed, gray] (b.east) -- (c.west);
  \draw[->, thick, dashed, gray] (c.east) -- (d.west);
  \draw[->, thick, dashed, gray] (e.east) -- (f.west);
  \node[gray, font=\scriptsize] at (3.0,-0.6) {solid = left-child};
  \node[gray, font=\scriptsize] at (3.05,-1.1) {dashed = right-sibling};
\end{tikzpicture}
\end{center}

Read it as a conventional tree: $A$ is the root with three children $B, C, D$; $B$ in turn has two children $E, F$. But in memory, $A$ holds only one downward pointer (to $B$, its first child), and $B, C, D$ are linked left-to-right by sibling pointers. Likewise $E$ points sideways to $F$. To list $A$'s children you go down to $B$, then follow dashed arrows: $B \to C \to D \to \textsc{nil}$. A node with $500$ children still needs only its two pointers; the $500$ children form a sibling chain among themselves.

\begin{focusbox}
The big idea of this section: \textbf{any rooted tree, no matter how many children per node, can be stored with just two pointers per node} via left-child/right-sibling. Trees do not require fancy variable-length nodes. This same two-pointer node is exactly the binary-tree node from before --- so deep down, an arbitrary tree and a binary tree share one representation. That unity is why binary trees are studied so intensely: they are secretly general.
\end{focusbox}

\subsection{Other tree representations}

Two more variations are worth naming. If you know an upper bound $k$ on the number of children (a ``$k$-ary tree''), you can give each node an array $child[1..k]$ of child pointers --- direct and fast, at the cost of wasting slots on nodes with fewer than $k$ children. And for the special case of a \emph{complete} binary tree --- one filled level by level with no gaps --- you can drop pointers entirely and store the whole tree in a plain array, using arithmetic to find relatives: the children of the node at index $i$ live at indices $2i$ and $2i+1$, and its parent at $\floor{i/2}$. That arithmetic trick is the engine behind the \emph{heap}, which gets its own chapter.

\begin{historybox}
In 1956 at the RAND Corporation, Allen Newell, J.C.\ Shaw, and Herbert Simon were attempting to
build the Logic Theorist --- the earliest artificial intelligence program, designed to discover
automated proofs for theorems in Russell and Whitehead's \emph{Principia Mathematica}. They quickly
hit a brick wall with standard contiguous arrays: the branching search trees of symbolic logic and
chess move exploration expanded and contracted unpredictably. If they allocated fixed-size arrays,
memory constantly ran out or sat mostly empty. In response, they designed the Information Processing
Language (IPL) and invented the \emph{linked list}, liberating data from the physical tyranny of
contiguous memory by letting each item point to the next wherever space was free. Concurrently in
Munich, Friedrich L.\ Bauer and Klaus Samelson (1957) patented the stack principle (the
\emph{Kellerprinzip}, or cellar principle) for translating mathematical formulas. Soon after, Edsger
W.\ Dijkstra (1960--1961) invented the Shunting-Yard algorithm, using an operator stack and an
output queue to parse nested parenthetical expressions for the world's first Algol 60 compiler, while
early operating system architects designed circular ring buffers to connect asynchronous hardware
interrupts and I/O devices with CPU consumers.

Yet the evolution of elementary data structures sparked intense practical controversies. For
decades, computer science textbooks taught that linked lists were strictly superior to arrays for
dynamic collections because they offer $\Oh(1)$ insertion and deletion without moving neighbors.
But modern computer architecture turned that advice upside down! To traverse a linked list, a
processor must ``chase pointers'' across randomly scattered heap addresses, stalling the CPU on
repeated cache misses. An array, though requiring $\Oh(n)$ shifts in the worst case, sits in
contiguous memory where hardware prefetchers can load whole cache lines at once, making contiguous
vectors (like C++'s \lstinline|std::vector|) dramatically faster for almost all practical workloads.
Furthermore, the ubiquitous null pointer --- introduced by Tony Hoare in 1965 for ALGOL~W, which he
later dubbed his ``billion-dollar mistake'' --- plagued programmers with endless null-dereference
crashes and messy special cases, until sentinel nodes unified list boundaries into clean circular
rings.

Mastering arrays, stacks, queues, and linked lists unlocks the fundamental scaffolding of all
systems software: compilers cannot parse grammar without stacks, operating systems cannot schedule
tasks without FIFO queues, and memory allocators cannot manage free space without linked lists. But
their honest limits must guide your choices: lists sacrifice $\Theta(1)$ random access, and on
64-bit architectures, pointer overhead can double or triple memory consumption. In Chapter~11, we
fuse contiguous arrays with linked-list chaining to construct Hash Tables with $\Oh(1)$ expected
search; in Chapter~12, we branch linear pointer chains into Binary Search Trees to achieve
$\Oh(\log n)$ search and insertion; and in Chapter~22, we deploy queues and stacks to drive
Breadth-First and Depth-First Graph Search.
\end{historybox}

\section{Putting it together}
\label{ch10:summary}

Step back and notice the single thread running through everything. The whole chapter rests on \emph{two} primitive ideas:

\begin{enumerate}
  \item The \textbf{contiguous block} (the array), which gives $\Th(1)$ random access by index via address arithmetic, but pays $\Th(n)$ to insert or delete in the middle because neighbors have no slack.
  \item The \textbf{pointer-linked object} (the node), which gives $\Th(1)$ structural edits by rewiring a few arrows, but pays $\Th(n)$ to reach the $i$-th element because there is no arithmetic shortcut --- you must walk.
\end{enumerate}

Stacks, queues, and deques are arrays (or lists) under a discipline about \emph{which end} you may touch, and they are all $\Th(1)$ per operation. Linked lists are pure pointer webs. Trees are pointer webs with a parent-child shape, and even the bushiest tree collapses to two pointers per node. Every advanced structure later in the book --- hash tables, balanced search trees, heaps, graphs --- is some artful combination of these two primitives, chosen to make the operations \emph{you} care about cheap.

\begin{recoverybox}
Before moving on, make sure you can answer these without peeking. Why is array random access $\Th(1)$ but array middle-insertion $\Th(n)$? Why is linked-list insertion $\Th(1)$ but linked-list random access $\Th(n)$? What does a stack's LIFO discipline buy you, and a queue's FIFO? Why does a circular array stop a queue from crawling off the end? And how do you store a tree where a node might have a thousand children using only two pointers per node? If those five answers feel solid, you are ready for what comes next.
\end{recoverybox}

\section*{Exercises}
\addcontentsline{toc}{section}{Exercises}

\begin{exercise}[Address arithmetic]
\label{ch10:ex-address}
An array of $64$-bit integers ($s = 8$ cells per element) begins at base address $4096$ and uses zero-based indexing. (a) At what address does $A[0]$ live? $A[5]$? $A[5]$ if the array were instead one-based? (b) An element is found at address $4216$; what is its index (zero-based)? Show the arithmetic in both directions.
\end{exercise}

\begin{exercise}[Stack trace]
\label{ch10:ex-stacktrace}
Starting from an empty stack, perform: \textsc{Push}(5), \textsc{Push}(8), \textsc{Push}(2), \textsc{Pop}, \textsc{Push}(11), \textsc{Pop}, \textsc{Pop}, \textsc{Pop}. List the value returned by each \textsc{Pop}, in order, and state the contents of the stack after every operation. What goes wrong on the final operation, and what should the implementation do about it?
\end{exercise}

\begin{exercise}[Balanced brackets]
\label{ch10:ex-brackets}
Describe an algorithm using a single stack that decides whether a string of brackets such as \lstinline|([]{()})| is properly balanced (every opener has a matching closer, correctly nested). Trace your algorithm on \lstinline|([)]| and explain precisely where and why it reports an error. What is the running time on a string of length $n$?
\end{exercise}

\begin{exercise}[Circular queue]
\label{ch10:ex-circqueue}
Using a circular array of size $m = 5$ (one-based) with $head = tail = 1$, trace: \textsc{Enq}(1), \textsc{Enq}(2), \textsc{Enq}(3), \textsc{Deq}, \textsc{Deq}, \textsc{Enq}(4), \textsc{Enq}(5), \textsc{Enq}(6), \textsc{Enq}(7). Show the array contents and both indices after each step, and mark exactly where a wrap-around occurs. Then explain the ``full vs.\ empty'' ambiguity in your own words and propose a fix.
\end{exercise}

\begin{exercise}[Reverse a singly linked list]
\label{ch10:ex-reverse}
Write pseudocode that reverses a singly linked list \emph{in place}, returning the new head, using only a constant number of extra pointers (no second list, no recursion). Hint: walk the list maintaining three pointers --- \emph{prev}, \emph{curr}, \emph{next} --- and at each node flip its \emph{next} pointer to aim at \emph{prev}. Trace it on the list $9 \to 16 \to 4 \to \textsc{nil}$, drawing the pointers after each step. Argue that it runs in $\Th(n)$ time and $\Th(1)$ extra space.
\end{exercise}

\begin{exercise}[Two stacks make a queue]
\label{ch10:ex-twostacks}
Implement a queue using two stacks $S_{in}$ and $S_{out}$, supporting only \textsc{Push} and \textsc{Pop} on each. \textsc{Enqueue} pushes onto $S_{in}$. \textsc{Dequeue}: if $S_{out}$ is empty, pop \emph{everything} from $S_{in}$ and push it onto $S_{out}$, then pop from $S_{out}$. Trace \textsc{Enq}(1), \textsc{Enq}(2), \textsc{Enq}(3), \textsc{Deq}, \textsc{Enq}(4), \textsc{Deq}, \textsc{Deq} and verify that values come out in FIFO order $1,2,3,4$. Why does pouring from $S_{in}$ to $S_{out}$ produce the correct order? (You will analyze its amortized cost in a later chapter.)
\end{exercise}

\begin{exercise}[Doubly linked delete]
\label{ch10:ex-delete}
Given a doubly linked list with a sentinel, write the two-line \textsc{List-Delete} and argue it correctly handles deleting (a) the only node, (b) the first node, (c) the last node, (d) a middle node --- all without any \textsc{if} statements. Why does the sentinel make all four cases identical?
\end{exercise}

\begin{exercise}[Array vs.\ list, in practice]
\label{ch10:ex-vs}
For each scenario, decide whether a dynamic array or a doubly linked list is the better fit, and justify in one or two sentences: (a) storing the pixels of an image you will read millions of times but never resize; (b) a text editor's buffer where the user constantly inserts and deletes characters in the middle; (c) a to-do list where you mostly append at the end and occasionally read the $k$-th item; (d) the ``undo history'' of an editor.
\end{exercise}

\begin{exercise}[Left-child, right-sibling traversal]
\label{ch10:ex-lcrs}
A rooted tree is stored in left-child, right-sibling form: each node $x$ has $x.child$ (first child or \textsc{nil}) and $x.sibling$ (next sibling or \textsc{nil}). Write a recursive procedure \textsc{Count-Nodes}$(x)$ that returns the number of nodes in the subtree rooted at $x$ (counting $x$ itself). Then explain why the procedure, despite the tree being arbitrary-arity, only ever follows two kinds of pointer. What is its running time on a tree of $n$ nodes?
\end{exercise}

\begin{exercise}[Heap-as-array preview, challenge]
\label{ch10:ex-heaparray}
A complete binary tree is stored in a one-based array, where the node at index $i$ has children at $2i$ and $2i+1$ and parent at $\floor{i/2}$. (a) Draw the tree for the array $[16, 14, 10, 8, 7, 9, 3]$. (b) Which array index holds the parent of the node at index $6$? Its sibling? (c) If the array has $n$ elements, which indices are leaves? (Hint: a node is a leaf exactly when it has no left child, i.e.\ when $2i > n$.) Verify your answer on the $7$-element example.
\end{exercise}
